% Ad Atticum 2.1 (ad-atticum-02-01)
% Source: https://raw.githubusercontent.com/PerseusDL/canonical-latinLit/master/data/phi0474/phi057/phi0474.phi057.perseus-lat1.xml
% Fetched by scripts/fetch_latin.py

% Dateline (Perseus): Scr. Romae m. Iun. a. 694 (60) .

\ciceroLetterOpener{CICERO ATTICO salutem}

\ciceroSection{1} Kal. Iuniis eunti mihi Antium et gladiatores M. Metelli cupide relinquenti venit obviam tuus puer. is mihi litteras abs te et commentarium consulatus mei Graece scriptum reddidit. in quo laetatus sum me aliquanto ante de isdem rebus Graece item scriptum librum L. Cossinio ad te perferendum dedisse; nam si ego tuum ante legissem, furatum me abs te esse diceres. quamquam tua illa (legi enim libenter) horridula mihi atque incompta visa sunt, sed tamen erant ornata hoc ipso quod ornamenta neglexerant et, ut mulieres, ideo bene olere quia nihil olebant videbantur. meus autem liber totum Isocrati myrothecium atque omnis eius discipulorum arculas ac non nihil etiam Aristotelia pigmenta consumpsit. quem tu Corcyrae, ut mihi aliis litteris significas, strictim attigisti, post autem, ut arbitror, a Cossinio accepisti. quem tibi ego non essem ausus mittere nisi eum lente ac fastidiose probavissem.

\ciceroSection{2} quamquam ad me rescripsit iam Rhodo Posidonius se, nostrum illud ὑπόμνημα cum legeret, quod ego ad eum ut ornatius de isdem rebus scriberet miseram, non modo non excitatum esse ad scribendum sed etiam plane deterritum. quid quaeris? conturbavi Graecam nationem. ita vulgo qui instabant ut darem sibi quod ornarent iam exhibere mihi molestiam destiterunt. tu , si tibi placuerit liber, curabis ut et Athenis sit et in ceteris oppidis Graeciae; videtur enim posse aliquid nostris rebus lucis adferre.

\ciceroSection{3} oratiunculas autem et quas postulas et pluris etiam mittam, quoniam quidem ea quae nos scribimus adulescentulorum studiis excitati te etiam delectant. fuit enim mihi commodum, quod in eis orationibus quae Philippicae nominantur enituerat civis ille tuus Demosthenes, et quod se ab hoc refractariolo iudiciali dicendi genere abiunxerat ut σεμνότερόσ τισ et πολιτικώτεροσ videretur, curare ut meae quoque essent orationes quae consulares nominarentur. quarum una est in senatu Kal. Ianuariis, altera ad populum de lege agraria, tertia de Othone, quarta pro Rabirio, quinta de proscriptorum filiis, sexta cum provinciam in contione deposui, septima quom Catilinam emisi, octava quam habui ad populum postridie quam Catilina profugit, nona in contione quo die Allobroges indicarunt, decima in senatu Nonis Decembribus. sunt praeterea duae breves, quasi ἀποσπασμάτια legis agrariae. hoc totum σῶμα curabo ut habeas; et quoniam te cum scripta tum res meae delectant, isdem ex libris perspicies et quae gesserim et quae dixerim; aut ne poposcisses; ego enim tibi me non offerebam.

\ciceroSection{4} quod quaeris quid sit quo te arcessam, ac simul impeditum te negotiis esse significas neque recusas quin, non modo si opus sit sed etiam si velim, accurras, nihil sane est necesse, verum tamen videbare mihi tempora peregrinationis commodius posse discribere. nimis abes diu, praesertim cum sis in propinquis locis, neque nos te fruimur et tu nobis cares. ac nunc quidem otium est, sed, si paulo plus furor Pulchelli progredi posset, valde ego te istim excitarem. verum praeclare Metellus impedit et impediet. quid quaeris? est consul φιλόπατρισ et, ut semper iudicavi, natura bonus.

\ciceroSection{5} ille autem non simulat sed plane tribunus pl. fieri cupit. qua de re quom in senatu ageretur, fregi hominem et inconstantiam eius reprehendi qui Romae tribunatum pl. peteret cum in Sicilia hereditatem se petere dictitasset, neque magno opere dixi esse nobis laborandum, quod nihilo magis ei liciturum esset plebeio rem publicam perdere quam similibus eius me consule patriciis esset licitum. iam cum se ille septimo die venisse a freto neque sibi obviam quemquam prodire potuisse et noctu se introisse dixisset in eoque se in contione iactasset, nihil ei novi dixi accidisse. ex Sicilia septimo die Romam; ante tribus horis Roma Interamnam. noctu introisti; idem ante. non est itum obviam; ne tum quidem quom iri maxime debuit. quid quaeris? hominem petulantem modestum reddo non solum perpetua gravitate orationis sed etiam hoc genere dictorum. itaque iam familiariter cum ipso cavillor ac iocor; quin etiam cum candidatum deduceremus, quaerit ex me num consuessem Siculis locum gladiatoribus dare. negavi . at ego inquit novus patronus instituam; sed soror, quae tantum habeat consularis loci, unum mihi solum pedem dat noli , inquam de uno pede sororis queri; licet etiam alterum tollas. non consulare inquies dictum. fateor ; sed ego illam odi male consularem. ea est enim seditiosa, ea cum viro bellum gerit neque solum cum Metello sed etiam cum Fabio, quod eos †esse in hoc esse† moleste fert.

\ciceroSection{6} quod de agraria lege quaeris, sane iam videtur refrixisse. quod me quodam modo molli bracchio de Pompei familiaritate obiurgas, nolim ita existimes, me mei praesidi causa cum illo coniunctum esse, sed ita res erat instituta ut, si inter nos esset aliqua forte dissensio, maximas in re publica discordias versari esset necesse. quod a me ita praecautum atque provisum est non ut ego de optima illa mea ratione decederem sed ut ille esset melior et aliquid de populari levitate deponeret. quem de meis rebus, in quas eum multi incitarant, multo scito gloriosius quam de suis praedicare; sibi enim bene gestae, mihi conservatae rei publicae dat testimonium. hoc facere illum mihi quam prosit nescio; rei publicae certe prodest. quid ? si etiam Caesarem cuius nunc venti valde sunt secundi reddo meliorem, num tantum obsum rei publicae?

\ciceroSection{7} quin etiam si mihi nemo invideret, si omnes, ut erat aequum, faverent, tamen non minus esset probanda medicina quae sanaret vitiosas partis rei publicae quam quae exsecaret. nunc vero, quom equitatus ille quem ego in clivo Capitolino te signifero ac principe conlocaram senatum deseruerit, nostri autem principes digito se caelum putent attingere si mulli barbati in piscinis sint qui ad manum accedant, alia autem neglegant, nonne tibi satis prodesse videor si perficio ut nolint obesse qui possunt?

\ciceroSection{8} nam Catonem nostrum non tu amas plus quam ego; sed tamen ille optimo animo utens et summa fide nocet interdum rei publicae; dicit enim tamquam in Platonis πολιτείᾳ , non tamquam in Romuli faece sententiam. quid verius quam in iudicium venire qui ob rem iudicandam pecuniam acceperit? censuit hoc Cato, adsensit senatus; equites curiae bellum, non mihi; nam ego dissensi. quid impudentius publicanis renuntiantibus? fuit tamen retinendi ordinis causa facienda iactura. restitit et pervicit Cato. itaque nunc consule in carcere incluso, saepe item seditione commota, aspiravit nemo eorum quorum ego concursu itemque ii consules qui post me fuerunt rem publicam defendere solebant. quid ergo? istos inquies mercede conductos habebimus? quid faciemus, si aliter non possumus? an libertinis atque etiam servis serviamus? sed , ut tu ais, ἅλισ σπουδῆσ .

\ciceroSection{9} Favonius meam tribum tulit honestius quam suam, Luccei perdidit. accusavit Nasicam inhoneste ac modeste tamen; dixit ita ut Rhodi videretur molis potius quam Moloni operam dedisse. mihi quod defendissem leviter suscensuit. nunc tamen petit iterum rei publicae causa. Lucceius quid agat scribam ad te cum Caesarem videro, qui aderit biduo.

\ciceroSection{10} quod Sicyonii te laedunt, Catoni et eius aemulatori attribuis Servilio. quid ? ea plaga nonne ad multos bonos viros pertinet? sed si ita placuit, laudemus, deinde in dissensionibus soli relinquamur?

\ciceroSection{11} Amalthea mea te exspectat et indiget tui. Tusculanum et Pompeianum valde me delectant, nisi quod me, illum ipsum vindicem aeris alieni, aere non Corinthio sed hoc circumforaneo obruerunt. in Gallia speramus esse otium. Prognostica mea cum oratiunculis prope diem exspecta et tamen quid cogites de adventu tuo scribe ad nos. nam mihi Pomponia nuntiari iussit te mense Quintili Romae fore. id a tuis litteris quas ad me de censu tuo miseras discrepabat.

\ciceroSection{12} Paetus, ut antea ad te scripsi, omnis libros quos frater suus reliquisset mihi donavit. hoc illius munus in tua diligentia positum est. si me amas, cura ut conserventur et ad me perferantur; hoc mihi nihil potest esse gratius. et cum Graecos tum vero diligenter Latinos ut conserves velim. tuum esse hoc munusculum putabo. ad Octavium dedi litteras; cum ipso nihil eram locutus; neque enim ista tua negotia provincialia esse putabam neque te in tocullionibus habebam. sed scripsi, ut debui, diligenter.
